\documentclass[a4paper]{article}
% Español (México) con XeLaTeX: fontspec para fuentes Unicode, polyglossia
% para la división silábica y los nombres del documento en la variante mexicana.
\usepackage{fontspec}
\usepackage{polyglossia}
\setdefaultlanguage[variant=mexican]{spanish}
% Los nombres chinos van como «pinyin (original)»: xeCJK da a los caracteres
% Han su propia fuente; sin ella XeLaTeX los omite con un aviso.
\usepackage{xeCJK}
\setCJKmainfont{FandolSong-Regular.otf}
% polyglossia no traduce los nombres de listings.
\AtBeginDocument{\providecommand{\lstlistingname}{}\renewcommand{\lstlistingname}{Listado}%
  \providecommand{\lstlistlistingname}{}\renewcommand{\lstlistlistingname}{Índice de listados}}
\usepackage{amsmath, amssymb}
\usepackage{graphicx}
\usepackage[margin=2.5cm]{geometry}
\usepackage[most]{tcolorbox}
\usepackage{etoolbox}
\usepackage{listings}
\usepackage{hyperref}
\usepackage{booktabs}
\usepackage{float}
\newtcolorbox{knowledgebox}[1]{enhanced, colback=blue!5!white, colframe=blue!75!black, colbacktitle=blue!75!black, coltitle=white, fonttitle=\bfseries, title={#1}, attach boxed title to top left={yshift=-2mm, xshift=2mm}, boxrule=1pt, sharp corners}
\newtcolorbox{importantbox}[1]{enhanced, colback=yellow!10!white, colframe=yellow!80!black, colbacktitle=yellow!80!black, coltitle=black, fonttitle=\bfseries, title={#1}, sharp corners}
\newtcolorbox{warningbox}[1]{enhanced, colback=red!5!white, colframe=red!75!black, colbacktitle=red!75!black, coltitle=white, fonttitle=\bfseries, title={#1}, sharp corners}
\lstset{language=Python, basicstyle=\ttfamily\small, keywordstyle=\color{blue}, stringstyle=\color{red!60!black}, commentstyle=\color{green!60!black}, breaklines=true, frame=single, numbers=left, numberstyle=\tiny\color{gray}, captionpos=b, extendedchars=true}
\newcommand{\notetitle}{Advantage Estimator: GRPO, RLOO, REINFORCE++}
\newcommand{\noteauthors}{Elaboradas a partir de materiales públicos del curso}
\newcommand{\notedate}{2025}
\newcommand{\videochannel}{Wudaokou Nashi (五道口纳什)}
\newcommand{\videopublishdate}{2025}
\newcommand{\videoduration}{aproximadamente 25 minutos}
\newcommand{\videourl}{}
\newcommand{\videocoverpath}{cover.jpg}
\newcommand{\repourl}{https://github.com/hqhq1025/ai-course-notes}

% Footer with repo info
\usepackage{fancyhdr}
\pagestyle{fancy}
\fancyhf{}
\fancyfoot[C]{\thepage}
\fancyfoot[R]{\footnotesize\color{gray}\href{\repourl}{AI Course Notes}}
\renewcommand{\headrulewidth}{0pt}
\renewcommand{\footrulewidth}{0pt}

\begin{document}
\begin{titlepage}\centering
{\Large Notas del curso\par}\vspace{1.2cm}{\huge\bfseries \notetitle\par}\vspace{0.8cm}{\large \noteauthors\par}\vspace{0.3cm}{\large \notedate\par}\vspace{1.2cm}
\ifdefempty{\videocoverpath}{}{\includegraphics[width=0.82\textwidth,height=0.45\textheight,keepaspectratio]{\videocoverpath}\par}
\vfill
\begin{tcolorbox}[width=0.9\textwidth, colback=black!2!white, colframe=black!60, sharp corners]
\textbf{Autor o canal del video}: \videochannel\par\textbf{Fecha de publicación}: \videopublishdate\par\textbf{Duración del video}: \videoduration\par
\ifdefempty{\videourl}{}{\textbf{Enlace del video}: \href{\videourl}{\nolinkurl{\videourl}}\par}\par
\textbf{Página del proyecto}: \href{\repourl}{\nolinkurl{\repourl}}\par
\end{tcolorbox}
\end{titlepage}
\tableofcontents\newpage

\section{Introducción}
En esta entrega se presentan las implementaciones de los distintos Advantage Estimator de los algoritmos de aprendizaje por refuerzo en veRL. Código central: \texttt{trainer/ppo/core\_algorithms.py}.

\section{Conceptos clave}
\subsection{Token Level vs. Sequence Level}
\begin{importantbox}{Dos granularidades}
\begin{itemize}
  \item \textbf{Token Level}: cada token tiene un valor de advantage independiente (como GAE)
  \item \textbf{Sequence Level}: toda la secuencia comparte un valor de advantage (como GRPO)
\end{itemize}
\end{importantbox}

\subsection{Lote vs. Group}
Lote es el mini-batch del entrenamiento. Group es, en GRPO, la agrupación de varios responses para un mismo prompt.

\section{Explicación detallada de cada algoritmo}
\subsection{GAE (Generalized Advantage Estimation)}
Estimación de advantage a nivel de token, que requiere un value model.

\subsection{GRPO Advantage}
A nivel de secuencia, con estandarización dentro del grupo, sin necesidad de un value model.

\subsection{REINFORCE / REINFORCE++ / RLOO}
Distintas estrategias de sustracción de baseline, para reducir la varianza.

\subsection{Resumen de la sección}
Los distintos estimadores de advantage implican compromisos entre la granularidad, la necesidad de un value model y el control de la varianza.

\newpage
\section{Elección del método: no hay bala de plata para la estimación de la ventaja}
GAE, GRPO y RLOO parecen diferir solo en la fórmula, pero su verdadera diferencia está en la granularidad de la señal de entrenamiento, la forma de controlar la varianza y la dependencia de módulos adicionales. Entender estas diferencias importa más que memorizar las siglas.

\begin{knowledgebox}{Tres cosas que hay que observar al elegir un Advantage Estimator}
\begin{itemize}
  \item Si se necesita una señal a nivel de token o a nivel de secuencia
  \item Si se está dispuesto a introducir un value model para aumentar la complejidad
  \item Si la tarea actual teme más una varianza grande o un sesgo de estimación grande
\end{itemize}
\end{knowledgebox}

\subsection{Resumen de la sección}
La elección del método de estimación de la ventaja es, en esencia, una decisión de compromiso entre granularidad, complejidad y varianza. Entender este compromiso resulta más útil que perseguir algoritmos nuevos sin criterio.

\newpage
\section{Síntesis y ampliación}
\begin{enumerate}
  \item GAE necesita un value model, GRPO no lo necesita
  \item GRPO logra la estimación de advantage mediante la normalización dentro del grupo
  \item RLOO usa leave-one-out como baseline
\end{enumerate}
\end{document}

